\chapter{Глобальная циклическая сшивка четырёх проходов}
\label{ch:v6-single-thread-global-cycle-sewing-gate}

\section{Постановка}

Локальный двойственный каркас показал, что одни и те же проходы могут
воспроизводить $Q_*$ и $T_*$. Теперь требуется более сильное утверждение:
проходы должны быть участками одной неразрывной нити, а не несколькими
замкнутыми волокнами.

Микроскопическое отсутствие ветвления означает, что каждый входящий
полупроход соединяется ровно с одним исходящим. На четырёх метках такое
правило является перестановкой $\sigma\in S_4$. Равенство входной и
выходной степеней обеспечивает локальное сохранение числа проходов, но
одну глобальную нить даёт только перестановка с единственным циклом длины
четыре. Это дискретная версия эйлеровой постановки
\cite{single-thread-euler1741}.

\section{Полный перебор четырёх меток}

В группе $S_4$ имеются $24$ перестановки со следующими типами циклов:
\begin{center}
\begin{tabular}{c|ccccc}
тип & $1+1+1+1$ & $2+1+1$ & $2+2$ & $3+1$ & $4$\\
\hline
число & $1$ & $6$ & $3$ & $8$ & $6$
\end{tabular}
\end{center}
Ровно шесть правил дают один четырёхцикл. Все они нечётны.

Правильная тетраэдральная группа текущего вакуума равна $A_4\subset
SO(3)$. В ней содержатся только тождество, восемь циклов типа $3+1$ и три
двойные транспозиции типа $2+2$. Поэтому
\begin{equation}
 \#\{\sigma\in A_4:\sigma\text{ является четырёхциклом}\}=0.
\end{equation}

\begin{theorem}[Одношаговый запрет правильной тетраэдральной сшивки]
Ни одна сохраняющая ориентацию симметрия правильного тетраэдрального
каркаса не сшивает четыре метки в один цикл. Максимальная орбита имеет
длину три.
\end{theorem}

\section{Упорядоченный вес усиливает запрет}

Рабочие доли проходов равны
\begin{equation}
 w=(p,q,q,q),
 \qquad p\ne q.
\end{equation}
Стабилизатор этого веса обязан фиксировать выбранную ось и может лишь
переставлять три поперечные оси. Он изоморфен $S_3$, содержит шесть
элементов и также не содержит четырёхцикла.

Канонический генератор остаточной $Z_3$ имеет тип
\begin{equation}
 (1)(234).
\end{equation}
Он сохраняет упорядоченный вес, но создаёт два независимых контура:
одиночный выбранный проход и трёхцикл поперечных проходов.

Это не означает, что одна нить несовместима с неравными временами
пребывания. Один четырёхцикл может иметь участки разной длины и тем самым
давать произвольные положительные доли $w_a$. Однако такое течение является
подвеской перестановки с отдельно выведенными временами пролёта, а не
стационарной перестановкой равных тактов.

\section{Четырёхтактный кандидат и его граница}

Проект уже содержит сдвиг порядка четыре
\begin{equation}
 c_4=(1234),
\end{equation}
который действительно является одним циклом. Но его действие на
тетраэдральном каркасе имеет
\begin{equation}
 \det O(c_4)=-1.
\end{equation}
Следовательно, он принадлежит полной группе $S_4$, но не правильной
вращательной группе $A_4\subset SO(3)$. Кроме того, при равных тактах он
не сохраняет рабочий вес: численный остаток равен
\begin{equation}
 \|P_{c_4}w-w\|=1.2278907848.
\end{equation}

Таким образом, ранняя четырёхтактная петля является точным кандидатом на
порядок посещения проходов, но проект ещё не вывел два необходимых моста:

\begin{enumerate}
 \item компенсацию нечётной пространственной ориентации существующей
 Real-структурой;
 \item времена пребывания на четырёх участках, пропорциональные $w_a$.
\end{enumerate}

\section{Real-удвоение не решает задачу автоматически}

Проверены три стандартных подъёма любой перестановки на восемь состояний
``ось $\times$ ориентация'': диагональный, с Real-переворотом и с
обращением стрелки. Ни один из $24$ подъёмов каждого типа не является
единственным восьмициклом.

Для $c_4$ получены типы $4+4$, $4+4$ и $2+2+2+2$. Для остаточного $Z_3$
получены $3+3+1+1$, $6+2$ и $2+2+2+2$. Поэтому простое произведение
перестановки с Real-инволюцией не сшивает две ориентации в одну нить.

\section{Почему коизометрия не является сшивкой}

Аффинная коизометрия $V$ имеет размер $3\times4$, ранг три и одномерное
ядро. Она канонически задаёт тетраэдральный каркас и его моменты, но не
является биекцией четырёх входов и четырёх выходов. Использовать $V$ как
физическое правило прохождения через узел означало бы смешать
реконструкцию наблюдаемого с микроскопической непрерывностью.

\begin{nogocriterion}
Равенство входной и выходной степеней, унитарность амплитуд или сохранение
топологического заряда сами по себе не доказывают одну нить. Требуется
единственная циклическая компонента и отдельный закон, запрещающий
пересечение нити через саму себя. Эквивалентность узлов сохраняется только
при допустимых локальных деформациях без разрезания и склейки
\cite{single-thread-reidemeister1927}.
\end{nogocriterion}

\section{Топологическая граница}

Класс $\pi_1(SO(3)/Z_3)=Z_6$ не позволяет произвольно оборвать вихревую
линию в упорядоченном объёме без выхода из вакуумного многообразия. Но он
не запрещает пересоединение линий, схлопывание замкнутой петли или
аннигиляцию противоположных зарядов. Более того, обменный мост Тома VI
явно показывает, что нейтральная пара может быть устранена в расширенном
пространстве конфигураций. Поэтому дефектная топология не заменяет закон
неразрывности фундаментальной нити.

\section{Следующий различающий тест}

Текущий родитель не выводит одну глобальную нить, но гипотеза не
опровергнута: в проекте уже имеется единственный четырёхцикл $c_4$.
Следующий гейт должен проверить его подвеску над четырьмя
тетраэдральными проходами.

Нулевая гипотеза: корневая фаза порядка четыре задаёт порядок посещения,
Real-структура компенсирует нечётную ориентацию, а оператор состояния $R$
канонически задаёт времена пребывания $\tau_a\propto w_a$ без нового
параметра.

Стоп-критерий: если назначение фазовых состояний тетраэдральным осям или
времён пребывания остаётся произвольным, одна глобальная нить не выведена.

\begin{protocol}
\begin{enumerate}
 \item локальное правило без ветвления: \textbf{перестановка};
 \item четырёхциклы в $S_4$: \textbf{шесть};
 \item четырёхциклы в $A_4$: \textbf{ноль};
 \item остаточный $Z_3$: \textbf{два контура $3+1$};
 \item стабилизатор рабочего веса: \textbf{фиксирует выбранную ось};
 \item проектный $C_4$: \textbf{один четырёхцикл};
 \item $C_4$ уже является допустимой сшивкой: \textbf{нет};
 \item стандартные Real-подъёмы дают один восьмицикл: \textbf{нет};
 \item коизометрия является биекцией проходов: \textbf{нет};
 \item одна глобальная нить опровергнута: \textbf{нет};
 \item одна глобальная нить выведена: \textbf{нет};
 \item следующий тест: \textbf{$C_4$-подвеска с временами пребывания}.
\end{enumerate}
\end{protocol}

Машинный сертификат сохранён в
\path{s2t_v6_single_thread_global_cycle_sewing_gate_results.json}.

\begin{thebibliography}{9}
\bibitem{single-thread-euler1741}
L.~Euler,
\emph{Solutio problematis ad geometriam situs pertinentis},
Commentarii Academiae Scientiarum Imperialis Petropolitanae 8 (1741),
128--140.

\bibitem{single-thread-birkhoff1946}
G.~Birkhoff,
\emph{Tres observaciones sobre el algebra lineal},
Universidad Nacional de Tucum\'an, Revista, Serie A 5 (1946), 147--151.

\bibitem{single-thread-reidemeister1927}
K.~Reidemeister,
\emph{Elementare Begr\"undung der Knotentheorie},
Abhandlungen aus dem Mathematischen Seminar der Universit\"at Hamburg 5
(1927), 24--32.
\end{thebibliography}